\chapter{Unit 5: Testing, Validation and Customer Experience}
\label{chap:unit5_rev}

\begin{revbox}[Unit 5 Essential Summary]
\begin{itemize}[leftmargin=1.2em, itemsep=2pt]
    \item \textbf{Usability Testing:} Evaluates if users can use the interface effectively and efficiently (*``Did we build it right?''*).
    \item \textbf{Design Validation:} Evaluates if the product solves the real core human problem (*``Did we build the right thing?''*).
    \item \textbf{``Show, Don't Tell'' Doctrine:} Put the prototype in front of the user; observe natural behavior without pitching or explaining.
\end{itemize}
\end{revbox}

\section{Testing Conduct \& Feedback Capture Matrix}

\begin{itemize}[leftmargin=1.2em, itemsep=3pt]
    \item \textbf{Neutral Facilitation:} Reassure user (*``We are testing the product, not you''*); use non-leading neutral prompts (*``What did you expect to happen?''*); utilize the Think-Aloud Protocol.
    \item \textbf{Goal-Oriented Task Scenarios:} Provide realistic context, constraints, and end goals without revealing button names or navigation sequences.
    \item \textbf{Feedback Capture Matrix:}
    \begin{itemize}
        \item \textbf{Likes (+):} Elements users praised or found intuitive.
        \item \textbf{Wishes ($\Delta$):} Constructive criticisms and unmet expectations.
        \item \textbf{Questions (?):} Confusions and queries raised during testing.
        \item \textbf{Ideas (!):} New feature suggestions sparked during sessions.
    \end{itemize}
\end{itemize}

\section{Usability vs. Customer Experience (CX) \& Journey Mapping}

\begin{table}[h!]
\centering
\small
\renewcommand{\arraystretch}{1.3}
\begin{tabularx}{\linewidth}{p{2.8cm}X X}
\toprule
\textbf{Dimension} & \textbf{Usability (UI/UX)} & \textbf{Customer Experience (CX)} \\
\midrule
\textbf{Scope} & Interface-specific task interaction. & Holistic end-to-end relationship across all channels. \\
\textbf{Timeframe} & Minutes of active product usage. & Entire customer lifecycle (discovery to offboarding). \\
\textbf{Touchpoints} & Screen buttons, menus, mobile app. & Human (couriers), Digital (apps), Physical (packaging). \\
\textbf{Key Metrics} & Task Completion Rate, Time on Task, SEQ. & Net Promoter Score (NPS), CSAT, Retention Rate. \\
\bottomrule
\end{tabularx}
\caption{Usability vs. Customer Experience (CX) Comparison.}
\end{table}

\section{Digital Accessibility: WCAG POUR Principles}

\begin{formulabox}[The 4 WCAG Principles]
\begin{enumerate}[leftmargin=1.2em, itemsep=2pt]
    \item \textbf{Perceivable (P):} Provide text alternatives (\texttt{alt}-text) for non-text media; maintain minimum contrast ratio: \textbf{4.5:1} for normal body text, \textbf{3:1} for large text.
    \item \textbf{Operable (O):} Full keyboard navigation without requiring mouse gestures; no keyboard traps.
    \item \textbf{Understandable (U):} Predictable navigation, inline form validation, and clear error recovery.
    \item \textbf{Robust (R):} Compatibility across browsers and assistive screen readers.
\end{enumerate}
\end{formulabox}
